\documentclass[a4paper,10pt]{article}
\usepackage[francais]{babel}
\usepackage{url}

\selectlanguage{francais}
\input{letter-icube-uds.sty}

\selectlanguage{francais}
\begin{document}


\name{St\'ephane \textsc{Genaud}}
\fonction{Ma\^itre de conf\'erences (HDR)}
\telephoneillkirch{+33 (0)3688 54542}
\email{genaud@unistra.fr}

\begin{letter}[Recommendation de Jey Nadarasa pour un doctorat sur le sujet\\%
{\it Modeling of soft tissues under Cryoablation and High Intensity Focused Ultrasound with the effect of blood flow.}%
]{à}{\'Ecole Doctorale MSII,\\Université de Strasbourg}

\vspace{-13mm}
Madame, Monsieur

J'ai eu le plaisir d'accueillir Jey Nadarasa pour un stage de Master 1 pendant
dix semaines en juin et juillet 2012. Ce stage s'inscrivait dans son cursus de
Master \textit{Calcul Scientifique et Sécurité Informatique} de l'UFR de 
mathématiques-informatique de l'Université de Strasbourg.

Son sujet était l'optimisation de parties précises d'un code de physique destiné
à la simulation de la fusion des plasmas aux bords d'un tokamak. Il était encadré
principalement dans ce travail par Matthieu Kuhn, doctorant issu de la même
formation de Master, dont le travail de thèse se déroule dans le cadre d'un projet
ANR sur ce même sujet.

Le travail de Jey a été la réécriture d'une méthode de discrétisation basé sur 
le schéma d'Arakawa afin d'améliorer les temps d'exécution. Les techniques employées
ont été diverses. D'abord, il a travaillé sur l'optimisation séquentielle du code en faisant de 
l'inlining ou en jouant sur les variables intermédiaires, puis en favorisant la 
localité spatiale des données en mémoire afin de mieux exploiter les caches matériels.
Ensuite, il a parallélisé les nids de boucle au coeur de la méthode afin 
de profiter d'architectures multi-cores. La parallélisation s'est faite en OpenMP.


Jey a montré de très grandes qualités d'adaptation pour s'approprier les notions
d'architecture et de parallélisation nécessaires à sa tâche. Du côté mathématiques,
il possèdait le bagage nécessaire pour manier avec aisance les expressions du schéma
numérique et les simplifier. La rapidité avec laquelle Jey s'est emparé du sujet est
remarquable. Il a fait également preuve de beaucoup d'autonomie, cherchant toujours
lui même les solutions avant de demander des conseils. Enfin, Jey est une personne
extrêmement ouverte et sympathique, toujours désireuse de partager ses idées avec
l'équipe.

En conclusion, je recommande vivement Jey pour une thèse de doctorat, à la 
fois pour ses qualités humaines, son sérieux et l'efficacité avec laquelle il pourrait
mener un travail de recherche dans le domaine de la simulation numérique. 


\end{letter}
\begin{flushright}
\includegraphics[width=.20\textwidth]{signgenaud.jpg}
\end{flushright}
%\newpage
%\bibliographystyle{plain}
%\bibliography{julien_blaise_ED}


\end{document}
